\subsection{Likelihood localization}
\label{sec:support-likelihood-localization}

The first step toward Theorem~\ref{thm:atom-count-limit} is a deterministic
comparison around a fit that already carries a small amount of mass away from
its central atom. The comparison shows that such edge mass cannot create KKT
contacts where the central directional field has a quantitative gap below
one.

It is convenient to work in score coordinates. Replace a component location
\(u\) by \(t=\sqrt{1-\eps_n}\,u\); this bijection preserves support size.
We reuse \(G\) for the transformed mixing measure and set
\[
    \kappa_n=\frac{\eps_n}{2(1-\eps_n)},
    \qquad
    \gamma_n=(1-\eps_n)^{-1}=1+2\kappa_n.
\]
If the central atom is at
\(\tau\), the likelihood ratio of the atom at \(\tau+a\) to the atom at
\(\tau\), evaluated at \(Z_i\), is
\[
    q_{n,\tau,a}(i)
    =\exp\left\{a(Z_i-\gamma_n\tau)
                    -\frac{\gamma_na^2}{2}\right\}.
\]
We also write
\[
    q_t(z)=\exp\left\{tz-\frac{\gamma_nt^2}{2}\right\}
\]
for the likelihood ratio of the score atom \(t\) to the atom at zero.
For a perturbation \(G=(1-W)\delta_\tau+G^{\rm out}\), write
\[
    m_G(i)
    =1-W+\int q_{n,\tau,u-\tau}(i)\,G^{\rm out}(du)
\]
for its likelihood ratio to \(\delta_\tau\), and define
\[
    D_{n,\tau}(a)=\frac1n\sum_{i=1}^nq_{n,\tau,a}(i),
    \qquad
    \Psi_G(\tau+a)=\frac1n\sum_{i=1}^n
        \frac{q_{n,\tau,a}(i)}{m_G(i)}.
\]
The KKT conditions are \(\Psi_G\le1\), with equality \(G\)-almost everywhere;
continuity of \(\Psi_G\) upgrades this to equality at every support point.

The deterministic no-propagation estimate below separates the total mass of
the edge atoms from the geometry of their locations.
\begin{lemma}
\label{lem:support-no-percolation}
For every measure \(G=(1-W)\delta_\tau+G^{\rm out}\) defined above,
\[
    \Psi_G(\tau+a)\le \frac{D_{n,\tau}(a)}{1-W},
    \qquad a\in\mathbb R.
\]
Consequently, if \(D_{n,\tau}\le1-\eta\) on a set \(A\) and
\(W\le\eta/2\), then \(G\) has no support point in \(\tau+A\).
\end{lemma}

\begin{proof}
Every likelihood ratio is nonnegative, so \(m_G(i)\ge1-W\). Substitution in
the definition of \(\Psi_G\) gives the first inequality, and
\((1-\eta)/(1-\eta/2)<1\) gives the second. \qedhere
\end{proof}

An NPMLE places almost all its mass near the true score origin. The following
coarse estimate isolates a central component before the edge analysis. For a
mixing measure \(G\) in score coordinates, put
\[
    m_G(z)=\int
       \exp\left\{tz-\frac{\gamma_nt^2}{2}\right\}G(dt),
    \qquad
    \Delta_n(G)=\int\{1-e^{-\kappa_n|t|^2}\}\,G(dt).
\]
If \(P\) is standard Gaussian measure, then
\(Pm_G=1-\Delta_n(G)\), and hence
\begin{equation}
\label{eq:support-population-defect}
    P\log m_G\le\log Pm_G
       =\log\{1-\Delta_n(G)\}\le-\Delta_n(G).
\end{equation}

The population loss yields a coarse uniform empirical bound. Its
polylogarithmic precision is sufficient to extract a central atom before the
critical-edge analysis begins.
\begin{lemma}
\label{lem:support-coarse-defect}
There is a constant \(C<\infty\) such that, uniformly over
\(0\le\eps_n<1\),
\[
    \Delta_n(\widehat G_n)
    =O_{\mathbb P}\left\{\frac{(\log n)^C}{n}\right\}
\]
for the NPMLE \(\widehat G_n\).
\end{lemma}

\begin{proof}
It suffices to optimize over mixing measures supported in
\(\operatorname{conv}\{Z_1/\gamma_n,\ldots,Z_n/\gamma_n\}\). Indeed, if
\(p\) is the Euclidean projection of \(t\) onto this convex hull, then
\[
 (t-p)Z_i-\frac{\gamma_n}{2}(t^2-p^2)
 =(t-p)(Z_i-\gamma_np)-\frac{\gamma_n}{2}(t-p)^2\le0,
\]
so moving \(t\) to \(p\) weakly increases every fitted density. With
probability tending to one this interval lies in
\([-C\sqrt{\log n},C\sqrt{\log n}]\).

For \(\mathcal M(R)=\{m_G:\supp(G)\subset[-R,R]\}\), we use the compact
Gaussian-mixture entropy bound from
\cite[Section~2]{supportDivergenceCompanion}; see also
\cite[Section~3]{ghosal2001entropies}:
\begin{equation}
\label{eq:support-coarse-entropy}
 \log N_{[]}\{\eta,\mathcal M(R),L^1(P)\}
 \le C(1+R)\{1+\log(1/\eta)\}^2,
 \qquad 0<\eta\le\frac12.
\end{equation}
The constant is uniform in \(\kappa_n\), since
\[
    m_G(z)\phi(z)
    =\int e^{-\kappa_n t^2}\phi(z-t)\,G(dt),
\]
and an \(L^1(P)\) bracket for \(m_G\) is an ordinary \(L^1\) bracket for
a Gaussian location mixture of total mass at most one. To pass from the cited
probability-mixture bound to subprobabilities, write the mixing measure as
\(s\widetilde G\). Replace the endpoints of the probability-mixture brackets
by their positive parts, which does not increase their width. Bracket
\(s\in[0,1]\) on an \(\eta/4\)-grid and use probability-mixture brackets of
width \(\eta/4\). If \(s_j\le s\le s_j+\eta/4\) and
\(l\le m_{\widetilde G}\le u\), then
\([s_jl,(s_j+\eta/4)u]\) has width at most
\(\eta/2+\eta^2/16<\eta\). The additional
\(O\{\log(1/\eta)\}\) entropy is absorbed by
\eqref{eq:support-coarse-entropy}. Bracket inequalities are understood
\(P\)-almost everywhere; the displayed classes are separable, so the suprema
below are measurable.

Fix a dyadic shell \(\delta<\Delta_n(G)\le2\delta\), and cover it by upper
brackets \(u\) of \(L^1(P)\)-width \(\delta/4\). If \(m_G\le u\), then
\(Pu\le1-3\delta/4\), and Markov's inequality gives
\[
 \mathbb P\{P_n\log m_G\ge-\delta/4,\ m_G\le u\}
 \le e^{n\delta/4}(Pu)^n\le e^{-n\delta/2}.
\]
For \(R=O(\{\log n\}^{1/2})\), the logarithm of the number of brackets is
at most \((\log n)^C\) whenever \(\delta\ge n^{-2}\). A union bound over
brackets and dyadic shells therefore shows that every measure with
\(\Delta_n(G)>C(\log n)^C/n\) has negative likelihood gain relative to
\(\delta_0\), with probability tending to one. Such a measure cannot be an
NPMLE, proving the claim.
\end{proof}

The defect bound does not rule out several nearby atoms. The local collapse
lemma shows that the NPMLE carries all central mass on one atom, with the
remaining edge mass smaller than the KKT gaps used later.
\begin{lemma}
\label{lem:support-fixed-core-collapse}
Assume \(\kappa_n=n^{-\alpha+o(1)}\) for a fixed
\(0<\alpha<1/4\). For a sufficiently small fixed \(r_0>0\), the NPMLE can,
with probability tending to one,
be written
\begin{equation}
\label{eq:support-central-decomposition}
    \widehat G_n=(1-W_n)\delta_{\tau_n}+G_n^{\rm out},
    \qquad
    \supp(G_n^{\rm out})\subset[-r_0,r_0]^c,
\end{equation}
where
\[
    W_n+|\tau_n|^2
    =O_{\mathbb P}\left\{\frac{(\log n)^C}{n\kappa_n}\right\}.
\]
Moreover, \(W_n=o_{\mathbb P}(\kappa_n)\).
\end{lemma}

\begin{proof}
Lemma~\ref{lem:support-coarse-defect} and
\(1-e^{-\kappa_n|t|^2}\ge c\kappa_n\min\{|t|^2,1\}\) imply the displayed
bounds for the mass outside \([-r_0,r_0]\) and for the second moment of the
restriction to that interval. Normalize the central restriction to a probability
measure \(G^{\rm cen}\), and let \(\tau\) and \(V\) be its mean and variance.
Cauchy--Schwarz gives the asserted bound for \(|\tau|^2\).

We show that replacing \(G^{\rm cen}\) by \(\delta_\tau\) strictly increases
the likelihood unless \(V=0\). Keep the outside part fixed and interpolate
\[
 G_u=(1-W)\{(1-u)\delta_\tau+uG^{\rm cen}\}+G^{\rm out},
 \qquad 0\le u\le1.
\]
By concavity it is enough to show that the derivative at \(u=0\) is negative.
If \(m_i^0\) is the likelihood ratio under \(G_0\), put
\(\omega_i=(1-W)q_\tau(Z_i)/m_i^0\). With \(h=t-\tau\), the derivative is
\begin{equation}
\label{eq:support-collapse-derivative}
 \sum_{i=1}^n\omega_i
 \int\left[
   \exp\left\{h(Z_i-\gamma_n\tau)
                   -\frac{\gamma_nh^2}{2}\right\}-1
 \right]G^{\rm cen}(dt).
\end{equation}

First replace \(\omega_i\) by one. The bound on \(\tau\) allows us to choose a
deterministic \(\delta_n\downarrow0\) such that
\(|\tau|\le\delta_n\sqrt{\kappa_n}\) with probability tending to one. On this
event, a compact empirical-process bound gives
\[
 \sup_{\substack{|\tau|\le\delta_n\sqrt{\kappa_n}\\|h|\le2r_0}}
 \left|
 \partial_h^2 P_n
 e^{h(Z-\gamma_n\tau)-\gamma_nh^2/2}
 -
 \partial_h^2 P
 e^{h(Z-\gamma_n\tau)-\gamma_nh^2/2}
 \right|
 =
 O_{\mathbb P}\{n^{-1/2}(\log n)^C\}
 =
 o_{\mathbb P}(\kappa_n).
\]
Indeed, after two differentiations the envelope on the fixed parameter
rectangle is a fixed polynomial in \(|Z|\) times \(e^{2r_0|Z|}\), and hence
has finite moments of every order. The population average is
\[
    F_\tau(h)=\exp\{-\kappa_nh^2-\gamma_n\tau h\}.
\]
Its second derivative is at most \(-\kappa_n\) on \([-2r_0,2r_0]\), for small enough
\(r_0\), with probability tending to one. Since
\(\int h\,G^{\rm cen}(dt)=0\), the unweighted version of
\eqref{eq:support-collapse-derivative} is at most
\(-cn\kappa_nV\).

It remains to restore the weights. Write
\[
    A_i=\int\left[
    e^{h(Z_i-\gamma_n\tau)-\gamma_nh^2/2}-1
    \right]G^{\rm cen}(dt),
    \qquad r_i^0=1-\omega_i.
\]
Thus the likelihood ratio of the NPMLE differs from \(m_i^0\) by the factor
\(1+\omega_iA_i\). Taylor's theorem and
\(\int h\,G^{\rm cen}(dt)=0\) give
\[
 \left|
 \int\left[
 e^{h(Z_i-\gamma_n\tau)-\gamma_nh^2/2}-1
 \right]G^{\rm cen}(dt)
 \right|
 \le
 C(1+Z_i^2)e^{2r_0|Z_i|+C}V.
\]
On an event of probability tending to one, \(\max_i|Z_i|\le C_0\sqrt{\log n}\)
for a fixed \(C_0\). Hence
\[
    \max_i|A_i|
    \le C(1+\log n)e^{2C_0r_0\sqrt{\log n}}V
    =o_{\mathbb P}(1).
\]

To bound \(\sum_i r_i^0\), let \(\widehat r_i\) be the posterior
responsibility of \(G^{\rm out}\) under the original NPMLE. Integrating its
KKT equality over \(G^{\rm out}\) gives the exact identity
\(\sum_i\widehat r_i=nW\). The denominator comparison yields, uniformly in
\(i\),
\[
    r_i^0=\widehat r_i(1+\omega_iA_i)
    \le \{1+o_{\mathbb P}(1)\}\widehat r_i.
\]
Consequently \(\sum_i r_i^0=O_{\mathbb P}(nW)\). Replacing \(\omega_i\) by
one in \eqref{eq:support-collapse-derivative} therefore changes it by at most
\[
    O_{\mathbb P}\!\left\{
      nW(1+\log n)e^{2C_0r_0\sqrt{\log n}}V
    \right\}
    =o_{\mathbb P}\{n\kappa_nV\}.
\]
Indeed, the ratio of the first quantity to \(n\kappa_nV\) is
\(O_{\mathbb P}(n^{-1+2\alpha+o(1)})\). The derivative is
therefore negative whenever \(V\ne0\). Optimality forces \(V=0\), which proves
\eqref{eq:support-central-decomposition}; the final assertion follows from
\(\alpha<1/4\).
\end{proof}


We now control the total number of KKT contacts. On a fixed score interval,
the likelihood converges to a Poisson optimization with an analytic KKT
function and hence finitely many zeros. At the moving lower endpoint, where
the variance penalty changes on the shorter scale \(u_n^{-2}\), a separate
zero-counting argument controls the crossover.

Fix for now \(0<\alpha<1/4\), assume
\(\eps_n=n^{-\alpha+o(1)}\), and put
\(\theta_\alpha=1-\sqrt\alpha\). Let \(N\) be the Poisson process from
Proposition~\ref{prop:poisson-edge-field}, with intensity
\(\mu(dx)=e^{-x}\,dx\), and let \(S\) be its compensated edge field. For a
compact interval \(K=[b,c]\Subset(1/2,1)\) and a finite positive measure
\(\nu\) on \(K\), define
\[
    F_\nu(x)=\int_K e^{\theta x}\,\nu(d\theta),
    \qquad
    h(y)=\log(1+y)-y,
\]
and
\begin{equation}
\label{eq:support-one-d-limit-functional}
    \mathcal J_N^K(\nu)
    =\int_KS(\theta)\,\nu(d\theta)
      +\int_{\mathbb R}h(F_\nu(x))\,N(dx).
\end{equation}
The last integral is almost surely finite. For \(x<0\), use
\(|h(y)|\le y^2/2\) and \(F_\nu(x)\le\nu(K)e^{bx}\). The resulting
envelope is integrable against \(N\) because \(2b>1\). On
\([0,\infty)\), the Poisson process has only finitely many points.

To obtain \eqref{eq:support-one-d-limit-functional}, place masses
\(w_j=s_j/\{nc_n(\theta_j)\}\) at finitely many score locations
\(\theta_j u_n\). If \(v_i\) is the resulting likelihood increment at the
\(i\)th observation, then
\[
    \sum_i\log(1+v_i)
    =\sum_iv_i+\sum_i\{\log(1+v_i)-v_i\}.
\]
The first sum converges to \(\sum_js_jS(\theta_j)\). At an edge observation
\(Z_i=u_n+x/u_n\), the second sum sees
\(v_i\to\sum_js_je^{\theta_jx}\), while ordinary observations contribute
only a negligible quadratic remainder. This gives the two terms in
\eqref{eq:support-one-d-limit-functional}.

\subsection{Compact Poisson likelihoods}

After localization, tightness of the atom count reduces to finiteness of the
active set in a compact limiting Poisson optimization. For a fixed realization
of the Poisson process, the next proposition gives this finiteness uniformly
over all maximizers.
\begin{proposition}
\label{prop:support-one-d-analytic-active-set}
The functional \(\mathcal J_N^K\) is concave and coercive on finite positive
measures on \(K\), and it has at least one maximizer. For every maximizer
\(\nu\), its first-variation field
\begin{equation}
\label{eq:support-one-d-limit-kkt}
    \Gamma_\nu(\theta)
    =S(\theta)-\int_{\mathbb R}
       e^{\theta x}\frac{F_\nu(x)}{1+F_\nu(x)}N(dx)
\end{equation}
is real analytic on \((1/2,1)\) and satisfies
\[
    \Gamma_\nu\le0\quad\hbox{on }K,
    \qquad
    \Gamma_\nu=0\quad\nu\hbox{-almost everywhere}.
\]
In particular, \(\supp(\nu)\) is finite. Almost surely, the support sizes are
bounded uniformly over all maximizers of \(\mathcal J_N^K\). More generally,
for every \(M<\infty\), the numbers of zeros of \(\Gamma_\nu\) on \(K\) are
bounded uniformly over all positive measures \(\nu\) with \(\nu(K)\le M\).
\end{proposition}

\begin{proof}
Concavity follows from concavity of \(\log(1+y)\). To prove coercivity, let
\(m=\nu(K)\to\infty\), choose \(k\) with \(2^k\le m<2^{k+1}\), and put
\(L_k=k\log2/(2c)\). On \(I_k=[-L_k-1,-L_k]\),
\[
    F_\nu(x)\ge me^{cx}\ge e^{-c}2^{k/2}.
\]
The count \(N(I_k)\) has mean \(\asymp2^{k/(2c)}\). Chernoff's bound and
Borel--Cantelli imply that it is at least half its mean for every large \(k\),
almost surely. Since \(h(y)\le-y/2\) for large \(y\), the contribution from
\(I_k\) is at most \(-C2^{k(1/2+1/(2c))}\). This dominates the linear term,
which is \(O(m)=O(2^k)\), because \(c<1\). Thus the functional is coercive.
On sets of bounded mass it is continuous for weak convergence: the negative
Poisson tail is uniform by \(|h(F_\nu(x))|\le C_Me^{2bx}\), and the positive
tail is finite. Compactness of bounded sets of positive measures on \(K\)
then gives existence.

Differentiation in the direction of a point mass gives
\eqref{eq:support-one-d-limit-kkt} and the KKT inequalities. Both terms extend
analytically to a complex neighborhood of each compact subinterval of
\((1/2,1)\); on the negative half-line the integrand in the second term is
bounded by \(Ce^{(b+\operatorname{Re}\theta)x}\), which is integrable when
\(b+\operatorname{Re}\theta>1\). The function \(\Gamma_\nu\) is not
identically zero, since \(S(\theta)\to-\infty\) as \(\theta\uparrow1\) and
the subtracted integral is nonnegative. Its zeros on \(K\) are therefore
finite, and the KKT conditions place \(\supp(\nu)\) among them.

For uniformity, let \(\nu_j(K)\le M\). Compactness of bounded sets of positive
measures on \(K\) gives a weakly convergent subsequence, along which the
functions \(\Gamma_{\nu_j}\) converge locally uniformly on a complex
neighborhood of \(K\). The limit is \(\Gamma_\nu\) for the limiting measure
and is not identically zero: \(S(\theta)\to-\infty\) as
\(\theta\uparrow1\), whereas the integral subtracted in
\eqref{eq:support-one-d-limit-kkt} is nonnegative. Choose a bounded
domain \(U\) containing \(K\), with closure in this neighborhood, such that
\(\Gamma_\nu\) has no zeros on \(\partial U\). Rouch\'e's theorem then bounds
the number of zeros in \(K\) for all sufficiently large \(j\). This proves
the uniform zero bound, which in turn bounds the support of every maximizer.
\end{proof}

To transfer the limiting problem to the finite-sample likelihood,
Lemma~\ref{lem:support-fixed-core-collapse} supplies a central atom
\(\tau_n\). Put \(Y_{n,i}=Z_i-\gamma_n\tau_n\). For a sign
\(\xi\in\{-1,1\}\) and \(\theta\in K\), the likelihood ratio relative to
that central atom is
\[
 q_{n,\xi,\theta}(Z_i)
 =\exp\left\{\xi\theta u_nY_{n,i}
       -\frac{\theta^2u_n^2}{2}-\lambda_n\theta^2\right\},
 \qquad \lambda_n=\kappa_nu_n^2.
\]
Given positive measures \(\nu_+,\nu_-\) on \(K\), place physical mass
\(\nu_\xi(d\theta)/\{nc_n(\theta)\}\) at
\(\tau_n+\xi\theta u_n\), removing the same total mass from the central atom.
Denote the resulting likelihood gain by
\(\mathcal L_{n,K}(\nu_+,\nu_-)\).
The deterministic part of a normalized first variation is
\begin{equation}
\label{eq:support-deterministic-penalty}
    \mathcal P_n(\theta)
    =\frac{1-e^{-\lambda_n\theta^2}}{c_n(\theta)}.
\end{equation}

Uniform convergence and a scaled-mass bound allow the optimizing measure,
rather than only a fixed perturbation, to pass to the limit.
\begin{proposition}
\label{prop:support-one-d-compact-likelihood}
If \(K\Subset(\theta_\alpha,1)\), then on every set
\(\nu_+(K)+\nu_-(K)\le M\),
\[
 \mathcal L_{n,K}(\nu_+,\nu_-)
 \Rightarrow
 \mathcal J_{N_+}^K(\nu_+)+\mathcal J_{N_-}^K(\nu_-)
\]
in the space of continuous functions on the compact set of such measure
pairs, where \(N_+\) and \(N_-\) are independent copies of \(N\). The first
variations converge jointly, locally uniformly in a fixed complex
neighborhood of every smaller interval \(K'\Subset K\).

If \(K\Subset(1/2,1)\) is allowed to contain \(\theta_\alpha\), an NPMLE
whose noncentral support lies in the two signed copies of \(K\) has total
scaled edge mass \(O_{\mathbb P}(1)\).

More precisely, let \(K=[b,c]\Subset(1/2,1)\) contain \(\theta_\alpha\) and
assume
\[
    \frac\alpha2+(1-b)^2<\frac12.
\]
For NPMLEs whose noncentral support lies in the signed copies of \(K\), let
\(\nu_{n,+},\nu_{n,-}\) be their scaled edge measures. The normalized
first-variation fields, after adding \(\mathcal P_n\), are jointly tight for
locally uniform convergence on a fixed complex neighborhood of every
\(K'\Subset K\). Along any subsequence on which
\(\nu_{n,\xi}\Rightarrow\nu_\xi\), their limits are the analytic fields
\(\Gamma_{\nu_\xi}\) in \eqref{eq:support-one-d-limit-kkt}.
\end{proposition}

\begin{proof}
For finitely atomic measures, write
\[
 v_i=\sum_{\xi=\pm1}\int_K
       \frac{q_{n,\xi,\theta}(Z_i)-1}{nc_n(\theta)}
       \,\nu_\xi(d\theta).
\]
The identity
\(\sum_i\log(1+v_i)=\sum_iv_i+
\sum_i\{\log(1+v_i)-v_i\}\) separates the two limits. Proposition
\ref{prop:poisson-edge-field} sends the linear term to
\(\int S_+d\nu_++\int S_-d\nu_-\). The deterministic damping vanishes
uniformly on \(K\Subset(\theta_\alpha,1)\). Translation by the central atom
is also negligible after normalization by \(c_n(\theta)\), because, for
\(b=\min K\),
\[
    \frac{u_n|\tau_n|}{c_n(b)}
    =O_{\mathbb P}\!\left\{(\log n)^C
      n^{-1/2+\alpha/2+(1-b)^2+o(1)}\right\}
    =o_{\mathbb P}(1).
\]
At a right-edge observation \(Z_i=u_n+x/u_n\), the positive part of \(v_i\)
converges to \(F_{\nu_+}(x)\), while the negative part is exponentially small;
the left edge is analogous. This yields the nonlinear Poisson integrals.

The convergence is uniform over bounded measure sets. Indeed, the physical
edge mass is at most \(2M/\{nc_n(b)\}\), so its square contributes \(o(1)\).
On a bounded edge window, the kernels and their first few derivatives are
uniformly bounded and equicontinuous. Below \(-A\), the likelihood remainder
is bounded by \(C_Me^{2bx}\), and its first variation by
\(C_Me^{(b+\operatorname{Re}\theta)x}\). The finite-sample analogues follow
from the exponential-moment bounds used in Proposition
\ref{prop:poisson-edge-field}. A finite weak net of the compact measure set,
followed by \(A\to\infty\), proves the asserted functional and analytic
convergence.

The bounded-window equicontinuity and the negative-tail envelopes
\(C_Me^{2bx}\) and \(C_Me^{(b+\operatorname{Re}\theta)x}\) also prove the
analytic tightness assertion for an interval that crosses
\(\theta_\alpha\). The random linear term still
converges to \(S\), while its deterministic mean is exactly
\(-\mathcal P_n\); adding \(\mathcal P_n\) removes the only term that is not
tight across the crossover. In the nonlinear term,
\(e^{-\lambda_n\theta^2}=1+o(1)\) uniformly on the fixed complex
neighborhood, so the envelopes \(C_Me^{2bx}\) and
\(C_Me^{(b+\operatorname{Re}\theta)x}\) on the negative tail remain valid. The
translation by \(\tau_n\) is negligible, together with all derivatives, by
the displayed condition on \(b\). Uniformity over bounded measure sets allows
the fields to be evaluated at the sample-dependent \(\nu_{n,\xi}\); after a
weakly convergent subsequence, their limits are precisely
\(\Gamma_{\nu_\xi}\).

For the mass bound, put \(m=\nu_+(K)+\nu_-(K)\) and choose a sign carrying at
least \(m/2\). Uniformly on \(K\), the linear term is at most
\(B_nm\), where \(B_n=O_{\mathbb P}(1)\); this remains true when \(K\)
slightly crosses \(\theta_\alpha\), because the normalized variance penalty
is nonpositive. Let \(N_n(I)\) denote the number of observations whose edge
coordinate lies in \(I\). For a suitable constant \(c_1>0\),
\[
    \lim_{L\to\infty}\liminf_{n\to\infty}
    \mathbb P\!\left\{N_n([-L-1,-L])\ge c_1e^L\right\}=1,
\]
by the edge-process convergence. Each such observation has likelihood
increment at least \((m/3)e^{-c(L+1)}\), where \(c=\max K<1\). Since
\(h(y)\le-c_2y\) for large \(y\), this block contributes at most
\(-c_3me^{(1-c)L}\); all other nonlinear terms are nonpositive. First
restrict to \(B_n\le B\), then choose \(L\) so that this negative coefficient
exceeds \(2B\), and finally take \(m\) large. The resulting likelihood gain
is negative, while zero edge mass has gain zero. Hence the maximizing scaled
mass is tight.
\end{proof}

\subsection{The moving lower edge}

The compact Poisson limit leaves one gap in the proof of support tightness: a
narrow crossover where the normalized variance penalty changes from dominant
to negligible. There the penalty varies exponentially faster than the random
analytic remainder, which yields the deterministic zero bound needed to
control KKT contacts at the moving lower edge.
\begin{lemma}
\label{lem:support-steep-crossing}
Let \(I\) be a compact interval, and suppose \(R_n\) and \(b_n\) are analytic
on a fixed complex neighborhood of \(I\) and real-valued on \(I\). Assume
that every subsequence has a further subsequence
converging locally uniformly to a nonzero analytic function. Let
\[
    \mathcal P_n(\theta)=\exp\{a_n\varphi(\theta)+b_n(\theta)\},
    \qquad a_n\to\infty,
\]
where \(\varphi\) is analytic on the same neighborhood, real-valued on \(I\),
with \(\varphi'<0\) there, and
\(\{b_n-b_n(\theta_0)\}\) is locally uniformly bounded there for one
\(\theta_0\in I\). Along every subsequence on which \(R_n\) converges, the
number of real zeros of \(R_n-\mathcal P_n\) on \(I\) is bounded.
\end{lemma}

\begin{proof}
Pass to a subsequence with \(R_n\to R\not\equiv0\). Away from small disjoint
complex disks around the finitely many zeros of \(R\), both \(R_n'/R_n\) and
\(b_n'\) are bounded. Hence, on each component where \(R_n>0\),
\[
 \frac{d}{d\theta}\log\frac{R_n(\theta)}{\mathcal P_n(\theta)}
 =\frac{R_n'(\theta)}{R_n(\theta)}-a_n\varphi'(\theta)-b_n'(\theta)>0
\]
for large \(n\). There is at most one intersection on that component and none
where \(R_n<0\).

It remains to count intersections near a zero \(\theta_0\) of \(R\), of
multiplicity \(m\). Analytic factorization and Rouch\'e's theorem give
\[
    R_n(z)=g_n(z)\prod_{j=1}^m(z-z_{n,j}),
\]
where \(z_{n,j}\to\theta_0\), while \(g_n\), \(1/g_n\), and
\(g_n'/g_n\) are uniformly bounded. On a real interval where \(R_n>0\), set
\(H_n=\log R_n-a_n\varphi-b_n\). Its critical points satisfy
\begin{equation}
\label{eq:support-steep-log-derivative}
 0=H_n'(\theta)=A_n(\theta)+\sum_{j=1}^m\frac1{\theta-z_{n,j}},
 \qquad
 A_n=\frac{g_n'}{g_n}-a_n\varphi'-b_n'.
\end{equation}
If \(c_0=\min_I(-\varphi')\), then \(A_n\ge c_0a_n/2\) on the real axis.
Thus every critical point lies within \(O(a_n^{-1})\) of one of the roots
\(z_{n,j}\).

In each cluster of such roots, rescale \(y=a_n(\theta-x_n)\). After a
subsequence, the rescaled roots either converge or leave every compact set.
Multiplying \(a_n^{-1}H_n'\) by the polynomial \(Q_n\) formed from the
convergent roots gives
\[
    F_n(y)=C_n(y)Q_n(y)+Q_n'(y),
    \qquad C_n\to C=-\varphi'(\theta_0)>0.
\]
The limit \(CQ+Q'\) is not identically zero: at a root of \(Q\) of
multiplicity \(q\), the derivative has order \(q-1\), whereas \(CQ\) has
order at least \(q\). Rouch\'e's theorem therefore bounds the number of
critical points in each cluster. Rolle's theorem then bounds the zeros of
\(H_n\), equivalently the intersections \(R_n=\mathcal P_n\). Summing over the
finitely many zero disks and complementary components proves the claim.
\end{proof}

We apply this lemma to the finite-sample KKT certificate. For the penalty
\(\mathcal P_n\) in \eqref{eq:support-deterministic-penalty}, on a fixed
neighborhood of \(\theta_\alpha\),
\[
    \frac{d}{d\theta}\log \mathcal P_n(\theta)
    =-(1-\theta)u_n^2+O(1),
\]
so its transition has width \(O(u_n^{-2})\). Choose
\(I=[b,c]\Subset(1/2,1)\) around \(\theta_\alpha\) with
\begin{equation}
\label{eq:support-crossover-translation-condition}
    \frac\alpha2+(1-b)^2<\frac12.
\end{equation}
For the fitted measure, normalize the KKT certificate by
\[
    D_{\tau_n,\xi}(\theta)
    =\frac1n\sum_{i=1}^nq_{n,\xi,\theta}(Z_i),
    \qquad
    \Gamma_{n,\xi}(\theta)
    =\frac{\Psi_{\widehat G_n}(\tau_n+\xi\theta u_n)-1}{c_n(\theta)}.
\]
Write the fitted likelihood ratio relative to its central atom as \(1+v_i\).
Since that atom has positive mass, its KKT equality gives
\(n^{-1}\sum_i(1+v_i)^{-1}=1\). Consequently,
\begin{align}
 \Gamma_{n,\xi}(\theta)
 &=\frac1{nc_n(\theta)}\sum_{i=1}^n
   \frac{q_{n,\xi,\theta}(Z_i)-1}{1+v_i}\notag\\
 &=\frac{D_{\tau_n,\xi}(\theta)-1}{c_n(\theta)}
 -\frac1{nc_n(\theta)}\sum_{i=1}^n
   \{q_{n,\xi,\theta}(Z_i)-1\}\frac{v_i}{1+v_i}.
\label{eq:support-finite-kkt-split}
\end{align}
There is no residual contribution from the \(-1\) in the last sum, because
\[
    \sum_i\frac{v_i}{1+v_i}
    =\sum_i\left(1-\frac1{1+v_i}\right)=0.
\]
Proposition~\ref{prop:poisson-edge-field} controls the first term. In the
second, an observation with right-edge coordinate \(x\) converges to
\(-e^{\theta x}F_{\nu_+}(x)/\{1+F_{\nu_+}(x)\}\), and the left edge is
analogous. The negative-tail envelope is
\(C_Me^{(b+\operatorname{Re}\theta)x}\), which is integrable because
\(b+\operatorname{Re}\theta>1\). Once
Lemma~\ref{lem:support-one-d-exclusion-margins} confines the noncentral
support to the signed copies of \(I\), the crossover assertion of
Proposition~\ref{prop:support-one-d-compact-likelihood} gives
\begin{equation}
\label{eq:support-moving-edge-decomposition}
    \Gamma_{n,\xi}(\theta)=R_{n,\xi}(\theta)-\mathcal P_n(\theta),
\end{equation}
where \(\{R_{n,\xi}\}\) is tight for locally uniform convergence on the
complex neighborhood, and every subsequential limit is the nonzero function
\(\Gamma_\nu\) in
\eqref{eq:support-one-d-limit-kkt}. The only extra term is the translation by
\(\tau_n\), which is negligible with all derivatives by
\eqref{eq:support-crossover-translation-condition} and
\[
    \sup_{\theta\in I}\frac{u_n|\tau_n|}{c_n(\theta)}
    =O_{\mathbb P}\!\left\{(\log n)^C
      n^{-1/2+\alpha/2+(1-b)^2+o(1)}\right\}
    =o_{\mathbb P}(1).
\]
Finally,
\[
    \mathcal P_n(\theta)
    =\exp\left\{u_n^2(-\theta+\theta^2/2)
       +\log n+\log(1-e^{-\lambda_n\theta^2})\right\}.
\]
On a fixed simply connected complex neighborhood of \(I\), bounded away from
zero, \(1-e^{-\lambda_nz^2}\) has no zeros for all large \(n\). We choose the
analytic logarithm there that is real on \(I\). After subtraction at one fixed point, this logarithm
and all its derivatives are locally uniformly bounded. Lemma
\ref{lem:support-steep-crossing} therefore bounds the number of contacts in
the moving crossover.

To exclude contacts outside a compact edge interval, we collect the
quantitative gaps already proved in the phase-diagram analysis. Put
\[
 \tau_{0,n}=\frac{\bar Z_n}{\gamma_n},
 \qquad
 D_{0,\xi}(\theta)=D_{n,\tau_{0,n}}(\xi\theta u_n),
\]
so \(D_{0,\xi}\) is the KKT field relative to the best one-atom fit.
\begin{lemma}
\label{lem:support-one-d-exclusion-margins}
Fix \(r_0>0\) and
\(0<\eta<\theta_\alpha-1/2\). There is \(c>0\) such that, with probability
tending to one,
\[
 \sup_{\xi=\pm1}\sup_{r_0/u_n\le\theta\le\theta_\alpha-\eta}
 D_{0,\xi}(\theta)\le1-c\kappa_n.
\]
Moreover, for every \(\delta>0\), there are \(\rho>0\) and \(a>0\), with
\(\rho^2<\alpha\), such that, with probability at least
\(1-\delta+o(1)\),
\[
 \sup_{\xi=\pm1}\sup_{1-\rho\le\theta<1}
 \frac{D_{0,\xi}(\theta)-1}{c_n(\theta)}\le-a,
 \qquad
 \sup_{\xi=\pm1}\sup_{\theta\ge1}D_{0,\xi}(\theta)\le0.9.
\]
\end{lemma}

\begin{proof}
For the lower range, choose
\(\theta_*\in(1/2,\theta_\alpha-\eta)\) with
\(\theta_*^2<1/2-\alpha\). The Gaussian-range estimate used in Lemma
\ref{lem:subthreshold-envelope} gives
\[
 \sup_{\xi=\pm1}\sup_{0<\theta\le\theta_*}
 \frac{|K_{n,\xi}(\theta)-1|}{\theta^2}
 =O_{\mathbb P}\{n^{-1/2+\theta_*^2+o(1)}\}
 =o_{\mathbb P}(\lambda_n).
\]
Since \(D_{0,\xi}=e^{-\lambda_n\theta^2}K_{n,\xi}\), this leaves a deficit
at least \(c\kappa_n\) for \(r_0/u_n\le\theta\le\theta_*\). On
\([\theta_*,\theta_\alpha-\eta]\), Proposition
\ref{prop:poisson-edge-field} and Lemma~\ref{lem:uniform-centering} make the
random error \(o_{\mathbb P}(\lambda_n)\), while the damping loss is at least
\(c\lambda_n\). This proves the first assertion.

For the upper endpoint, use the decomposition in the proof of
Lemma~\ref{lem:endpoint-exclusion}. Fix \(M\). After points above the edge
level \(M\) are removed, the compensated field and its derivative are tight
on \([3/4,1]\). Thus, for a suitable \(K=K(M,\delta)\), both truncated fields
have supremum at most \(K\) with probability at least \(1-\delta/4+o(1)\).
The probability that either edge process has a point above \(M\) is at most
\(2e^{-M}+o(1)\). On the complementary event, the omitted compensator is at
least
\[
    \frac{e^{-(1-\theta)M}}{1-\theta},
    \qquad 1-\rho\le\theta<1.
\]
Choose \(M\) so that \(2e^{-M}<\delta/4\), and then choose \(\rho>0\) so
small that this lower bound exceeds \(K+2\) throughout the interval. The
normalized null field is then at most \(-2\). The damping is nonpositive, and
\(\rho^2<\alpha\) makes the empirical-centering error \(o_{\mathbb P}(c_n)\);
after decreasing the margin, this proves the first upper-endpoint inequality
with \(a=1\).

For \(\theta\ge1\), choose \(R\) so large that both centered sample extremes
are below the score \(1+R/u_n^2\), except on an event of probability at most
\(\delta/4+o(1)\). The right-of-one estimate in the proof of
Lemma~\ref{lem:endpoint-exclusion}, applied up to this score, gives
\(D_{0,\xi}\le0.9\). Beyond it, Lemma~\ref{lem:one-atom-D} makes the field
monotone away from the corresponding sample extreme. This proves the second
inequality, simultaneously for both signs, with total failure probability
at most \(\delta+o(1)\).
\end{proof}

The limiting-law argument requires compactness of the finite-sample active
sets before the limiting optimizer can be identified. The central-atom
localization, the exclusion margins, and the analytic zero bound combine to
give this fixed-exponent tightness.
\begin{lemma}
\label{lem:fixed-exponent-support-tightness}
If \(r_n\to\alpha\in(0,1/4)\), then
\[
    |\supp(\widehat G_n)|=O_{\mathbb P}(1).
\]
\end{lemma}

\begin{proof}
Lemma
\ref{lem:support-fixed-core-collapse} gives one central atom and outside mass
\[
    W_n=O_{\mathbb P}\{(\log n)^C/(n\kappa_n)\}
        =o_{\mathbb P}(\kappa_n).
\]
Relative to the best one-atom center, write
\(D_{\tau_n,\xi}(\theta)=D_{n,\tau_n}(\xi\theta u_n)\). The central atom
introduces the exact translation
\[
 D_{\tau_n,\xi}(\theta)
 =e^{-\xi\theta u_n\delta_n}D_{0,\xi}(\theta),
 \qquad \delta_n=\gamma_n\tau_n-\bar Z_n,
\]
where
\[
    |\delta_n|
    =O_{\mathbb P}\!\left\{(\log n)^C
      n^{-1/2+\alpha/2+o(1)}\right\}.
\]

Choose \(\eta>0\) so small that
\(\alpha/2+(\sqrt\alpha+2\eta)^2<1/2\). The first bound in Lemma
\ref{lem:support-one-d-exclusion-margins}, the translation estimate, and
Lemma~\ref{lem:support-no-percolation} exclude contacts between the fixed core
and \((\theta_\alpha-\eta)u_n\). For the upper endpoint choose \(\rho\) from
the second bound, with \(1-\rho>\theta_\alpha-\eta\). There the normalized
gap is at least \(ac_n(1-\rho)=n^{-\rho^2+o(1)}\), whereas
\[
    W_n=O_{\mathbb P}\!\left\{(\log n)^C
      n^{-1+\alpha+o(1)}\right\},
    \qquad
    u_n|\delta_n|
    =O_{\mathbb P}\!\left\{(\log n)^C
      n^{-1/2+\alpha/2+o(1)}\right\}.
\]
Both are negligible relative to the gap because \(\rho^2<\alpha<1/4\).
It remains only to note that a support point cannot have an arbitrarily large
translated score. The convex-hull localization in the proof of
Lemma~\ref{lem:support-coarse-defect} places every score atom of
\(\widehat G_n\) in
\([\min_iZ_i/\gamma_n,\max_iZ_i/\gamma_n]\). After increasing the fixed
\(R\) in Lemma~\ref{lem:support-one-d-exclusion-margins}, this implies that
every displacement from \(\tau_n\) has
\[
    \theta\le1+\frac{R+1}{u_n^2}
\]
outside an event of probability \(\delta+o(1)\); here
\(u_n|\tau_n|=o_{\mathbb P}(1)\). On this bounded endpoint window the
translation factor \(e^{-\xi\theta u_n\delta_n}\) is uniformly
\(1+o_{\mathbb P}(1)\). The normalized gap for \(1-\rho\le\theta<1\), the
fixed (0.1) gap for \(1\le\theta\le1+(R+1)/u_n^2\), and
Lemma~\ref{lem:support-no-percolation} therefore exclude all contacts beyond
\((1-\rho)u_n\), apart from an event of probability \(O(\delta)+o(1)\).

Every noncentral support point now lies in a signed copy of
\([\theta_\alpha-\eta,1-\rho]\), which is compactly contained in
\[
    K=[\theta_\alpha-2\eta,1-\rho/2]\Subset(1/2,1).
\]
Proposition~\ref{prop:support-one-d-compact-likelihood} makes the scaled mass
on \(K\) tight and gives joint analytic convergence on a complex neighborhood
of the smaller support interval.
Prokhorov's theorem and a Skorokhod representation therefore allow us, from
every subsequence, to extract a further subsequence on which the edge measures
converge weakly and the analytic remainders in
\eqref{eq:support-moving-edge-decomposition} converge locally uniformly,
almost surely on a common probability space.
Away from a fixed neighborhood of \(\theta_\alpha\), the uniform analytic
zero bound in Proposition~\ref{prop:support-one-d-analytic-active-set} applies
to the limiting remainders. Inside that neighborhood,
Lemma~\ref{lem:support-steep-crossing} controls the intersections with the
rapidly varying penalty. Rouch\'e's theorem transfers both bounds to the exact
KKT certificates along the subsequence. Since every support point is a zero,
the support counts are tight. Otherwise one could choose a subsequence
on which they escape every deterministic bound with fixed positive
probability. Repeating the Prokhorov--Skorokhod subsequence extraction just
described would then give a contradiction.

\end{proof}
